\chapter{Extra: CMS e GraphQL}

\begin{errorbox}{{Contenuti extra, fuori programma}}
  Il contenuto di questo capitolo \textbf{non fa parte del programma} di
  Tecnologie Web: corrisponde a una lezione bonus tenuta dal docente. Questi
  argomenti \textbf{non compaiono nelle prove scritte} né vengono chiesti
  durante la discussione del progetto.

  Per quanto riguarda il progetto, si ricordi che la traccia
  \textbf{vieta esplicitamente} l'uso di CMS, mentre \textbf{ammette} l'uso di
  GraphQL.

  Il capitolo è incluso per completezza, ma può essere saltato senza
  conseguenze ai fini dell'esame.
\end{errorbox}

\section{Il problema di gestire i contenuti}

L'obiettivo principale di molte applicazioni web è \textbf{mostrare contenuti}
agli utenti. Questi contenuti possono cambiare molto di frequente e crescere in
modo esponenziale: Wikipedia ne conta circa sette milioni di pagine.

Da qui alcune difficoltà:

\begin{itemize}
  \item scrivere a mano intere pagine web per aggiornare i contenuti richiede
        competenze tecniche e rallenta il lavoro;
  \item chi si occupa dei contenuti tipicamente \textbf{non} è chi sviluppa;
  \item gestire versioni diverse del contenuto pubblicato non è banale;
  \item garantire l'adozione di un flusso di lavoro strutturato è complicato.
\end{itemize}

\dfn{Content Management System}{
  Un \term{CMS} è un'applicazione web progettata per consentire la modifica
  collaborativa, l'organizzazione e la pubblicazione di contenuti
  \textbf{senza scrivere una riga di codice}, usando l'interfaccia fornita dal
  CMS stesso. Permette a utenti non tecnici di gestire e aggiornare i contenuti
  senza dipendere continuamente dagli sviluppatori.}

I CMS sono spesso usati per blog, siti di notizie e siti di commercio
elettronico. Le funzionalità tipiche comprendono:

\begin{itemize}
  \item pubblicazione web dei contenuti;
  \item formattazione, tipicamente con un editor grafico WYSIWYG;
  \item controllo di versione dei contenuti;
  \item indicizzazione, ricerca e recupero;
  \item autenticazione e autorizzazione con un sistema di permessi per gruppi;
  \item gestore di contenuti multimediali integrato;
  \item estensibilità e personalizzazione tramite temi ed estensioni.
\end{itemize}

\section{Architetture}

\subsection{CMS tradizionali}

Un CMS tradizionale è composto da due parti:

\begin{itemize}
  \item la \term{Content Management Application} (CMA), che permette agli
        utenti non tecnici di gestire i contenuti;
  \item la \term{Content Delivery Application} (CDA), che rende i contenuti
        aggiornati disponibili agli utenti finali, per esempio come pagine web.
\end{itemize}

L'architettura è \textbf{monolitica}: CMA e CDA sono strettamente integrate. Il
vantaggio è la facilità di distribuzione e di gestione (una sola applicazione).
Gli svantaggi sono la scarsa flessibilità in alcuni scenari e la limitazione dei
canali di distribuzione: i CMS tradizionali privilegiano l'output web, e non è
immediato esporre una API REST.

\subsection{CMS headless e disaccoppiati}

\dfn{CMS headless}{
  Un \term{CMS headless} separa la CMA (il \textit{corpo}, che si occupa
  dell'organizzazione dei contenuti) dalla CDA (la \textit{testa}, che si occupa
  della presentazione). È un CMS \textbf{senza testa}: non gestisce la
  presentazione ed espone soltanto una API. In questo modo i contenuti possono
  essere distribuiti su qualunque canale.}

I \term{CMS disaccoppiati} sono CMS tradizionali in cui CMA e CDA sono
separate: si può usare la CDA inclusa oppure realizzare modalità di
distribuzione diverse tramite le API esposte dalla CMA. Oggi la maggior parte
dei CMS tradizionali è disaccoppiata, WordPress compreso.

\begin{figure}[H]
  \centering
  \begin{tikzpicture}[font=\Lato\scriptsize,
      b/.style={rounded corners=3pt, line width=0.8pt, align=center,
                minimum width=2.7cm, minimum height=0.6cm},
      db/.style={b, draw=neutral!55!white, fill=neutral!8!white},
      be/.style={b, draw=primary!65!black, fill=primary!10!white},
      fe/.style={b, draw=accent!60!black, fill=accent!10!white},
      api/.style={b, draw=green-gradient-6, fill=green-gradient-1!45!white},
      a/.style={-{Stealth[length=4pt,width=3.5pt]}, line width=0.7pt,
                draw=neutral!45!black}]

    \foreach \x/\ttl in {0/{Tradizionale}, 4.6/{Disaccoppiato}, 9.2/{Headless}} {
      \node[font=\Lato\scriptsize\bfseries, text=neutral!25!black]
        at (\x,2.5) {\ttl};
      \node[db] at (\x,1.9) {DBMS};
      \node[be] at (\x,1.1) {UI di back-end};
    }

    % tradizionale
    \node[fe] (f0) at (0,0.3) {front-end\\ (template)};
    \node[db, draw=neutral!45!white] at (0,-0.7) {HTML};
    \draw[a] (0,0.0) -- (0,-0.45);

    % disaccoppiato
    \node[fe] (f1) at (4.6,0.3) {front-end};
    \node[api] (a1) at (4.6,-0.6) {API};
    \node[db, draw=neutral!45!white] at (3.6,-1.5) {HTML};
    \node[db, draw=neutral!45!white] at (5.7,-1.5) {app};
    \draw[a] (4.0,-0.9) -- (3.7,-1.25);
    \draw[a] (5.2,-0.9) -- (5.6,-1.25);

    % headless
    \node[api] (a2) at (9.2,0.3) {API};
    \node[db, draw=neutral!45!white] at (8.2,-0.7) {web app};
    \node[db, draw=neutral!45!white] at (10.3,-0.7) {app mobile};
    \draw[a] (8.7,0.0) -- (8.3,-0.45);
    \draw[a] (9.7,0.0) -- (10.2,-0.45);
  \end{tikzpicture}
  \caption{Le tre architetture a confronto.}
  \label{fig:cms-arch}
\end{figure}

\error{I CMS sono bersagli privilegiati}{
  Lo stesso CMS è usato su un numero enormi di siti: una vulnerabilità in un CMS
  può essere sfruttata su moltissimi siti contemporaneamente. È quindi
  imperativo \textbf{tenerli aggiornati} per applicare le correzioni al più
  presto.

  Un esempio concreto: la vulnerabilità \code{CVE-2022-21661} rendeva WordPress
  vulnerabile a SQL injection nelle versioni dalla 3.7.37 alla 5.8.3.}

\section{CMS tradizionali}

Fra i CMS più noti: \term{WordPress} (il più diffuso, usato da circa il 43\% di
tutti i siti web), \term{Joomla} (che alimenta il sito dei corsi di Informatica
della Federico II), \term{WooCommerce} (per il commercio elettronico) e
\term{MediaWiki} (che alimenta Wikipedia, Fandom, WikiHow).

Un CMS tradizionale è un'applicazione web come le altre:

\begin{itemize}
  \item nella fase di installazione chiede i dati per connettersi a un database
        relazionale (URL, nome utente, password) e crea le tabelle necessarie;
  \item chiede le credenziali per l'account di amministrazione;
  \item dopo l'installazione, l'amministratore può creare e organizzare
        contenuti, aggiungere utenti, definire i template e personalizzare
        l'aspetto del sito.
\end{itemize}

\section{Strapi, un CMS headless}

Fra i CMS headless più diffusi: \term{Strapi} (basato su Node.js),
\term{Wagtail} (basato su Django, per Python), \term{Ghost} (Node.js) e
\term{CrafterCMS} (Java).

\begin{lstlisting}[style=shell]
npx create-strapi-app@latest project-name
cd project-name
npm run develop
\end{lstlisting}

\begin{console}[Terminale]
> todo-api@0.1.0 develop
> strapi develop

Building build context (577ms)
Creating admin (16224ms)
Loading Strapi (1505ms)
Generating types (1413ms)

Welcome back!
To manage your project, go to the administration panel at:
http://localhost:1337/admin
To access the server, go to:
http://localhost:1337
\end{console}

Dall'interfaccia di amministrazione si definisce la \guillemotleft forma\guillemotright\ delle
risorse (i loro campi e i loro tipi), si stabiliscono le \textbf{relazioni} fra
risorse con un editor visuale, e infine si creano i contenuti veri e propri.

\subsection{La API REST generata}

Per ogni risorsa, Strapi genera automaticamente gli endpoint REST.

\begin{center}
\renewcommand{\arraystretch}{1.4}
\begin{tabular}{@{}l@{\hspace{0.8cm}}l@{\hspace{0.8cm}}>{\raggedright\arraybackslash}p{5.0cm}@{}}
  \textbf{Metodo} & \textbf{URL} & \textbf{Descrizione} \\
  \code{GET}    & \code{/api/:pluralApiId}             & elenco delle voci \\
  \code{POST}   & \code{/api/:pluralApiId}             & crea una voce \\
  \code{GET}    & \code{/api/:pluralApiId/:resourceId} & una voce specifica \\
  \code{PUT}    & \code{/api/:pluralApiId/:resourceId} & aggiorna una voce \\
  \code{DELETE} & \code{/api/:pluralApiId/:resourceId} & elimina una voce \\
\end{tabular}
\end{center}

Dall'interfaccia si creano token di autorizzazione (\textit{Settings > API
Tokens}), specificando per ciascuno un nome, una durata e i permessi di lettura
e scrittura per ogni risorsa. Gli utenti devono passare il token nell'header
\code{Authorization}, esattamente come nello schema JWT visto nel capitolo 15.

\begin{lstlisting}[style=http, name=Richiesta]
GET http://localhost:1337/api/professors
Authorization: Bearer 97060c911...e12
\end{lstlisting}

\begin{lstlisting}[style=json, name=Risposta]
{
  "data": [
    {
      "id": 1,
      "attributes": {
        "firstName": "Luigi Libero Lucio",
        "lastName": "Starace",
        "email": "luigiliberolucio.starace@unina.it",
        "createdAt": "2023-12-29T07:19:19.746Z",
        "updatedAt": "2023-12-29T07:25:16.666Z",
        "publishedAt": "2023-12-29T07:20:35.307Z"
      }
    }
  ]
}
\end{lstlisting}

\subsection{Parametri di query}

Gli endpoint di Strapi supportano numerosi parametri nella query string per
filtrare, ordinare e paginare i risultati, selezionare campi e popolare le
relazioni.

\begin{center}
\renewcommand{\arraystretch}{1.4}
\begin{tabular}{@{}l@{\hspace{1.2cm}}>{\raggedright\arraybackslash}p{8.6cm}@{}}
  \textbf{Parametro} & \textbf{Effetto} \\
  \code{populate}   & quali risorse collegate includere nei risultati \\
  \code{fields}     & quali campi includere, e in quale ordine \\
  \code{filters}    & interrogazioni personalizzate per filtrare i dati \\
  \code{sort}       & come ordinare i risultati \\
  \code{pagination} & la paginazione dei risultati \\
\end{tabular}
\end{center}

\begin{lstlisting}[style=http, name=Esempi]
GET /api/professors?populate=courses
GET /api/professors?fields[0]=lastName&fields[1]=firstName

# corsi da 6 crediti
GET /api/courses?filters[credits][$eq]=6

# corsi da almeno 6 crediti il cui titolo contiene "Web"
GET /api/courses?filters[credits][$gte]=6&filters[title][$contains]=Web

# professori il cui nome contiene "Luigi" (senza distinzione di maiuscole)
GET /api/professors?filters[firstName][$containsi]=Luigi
\end{lstlisting}

\section{GraphQL}

\subsection{I limiti di REST}

REST è un modo pratico e diffusissimo di esporre dati da un server. In alcuni
scenari, però, può risultare inefficiente.

\begin{center}
\renewcommand{\arraystretch}{1.45}
\begin{tabular}{@{}l@{\hspace{1.0cm}}>{\raggedright\arraybackslash}p{8.8cm}@{}}
  \textbf{Problema} & \textbf{In che consiste} \\
  \term{Over-fetching}
    & talvolta non ci servono tutti i campi di una risorsa, ma li scarichiamo
      comunque \\
  \term{Under-fetching}
    & talvolta ci servono risorse collegate, e servono richieste aggiuntive per
      ottenerle: prima l'utente, poi i suoi post, poi per ogni post i commenti\ldots \\
  \textbf{Evoluzione dell'API}
    & le API cambiano nel tempo: vanno versionate, va mantenuta una certa
      retro-compatibilità, alcune funzionalità vanno deprecate \\
\end{tabular}
\end{center}

Entrambi i primi due degradano le prestazioni e aumentano il traffico: nel primo
caso si scaricano dati inutili, nel secondo servono molte richieste aggiuntive.

\dfn{GraphQL}{
  \term{GraphQL} è un linguaggio di interrogazione per API, sviluppato da
  Facebook e pubblicato in versione open source nel 2015. Invece di appoggiarsi
  a endpoint predefiniti come in REST, permette a chi consuma l'API di
  \textbf{specificare in modo dichiarativo quali dati gli servono}, tramite
  interrogazioni.}

\subsection{GraphQL con Strapi}

\begin{lstlisting}[style=shell]
npm install @strapi/plugin-graphql
\end{lstlisting}

Dopo il riavvio, un \textit{playground} GraphQL è disponibile su
\code{http://localhost:1337/graphql}.

\begin{esempio}{Solo i campi che servono}
Supponiamo di dover mostrare soltanto il cognome di ciascun professore.

\begin{lstlisting}[style=plain, name=query]
query {
  professors {
    data {
      attributes {
        lastName
      }
    }
  }
}
\end{lstlisting}

\begin{lstlisting}[style=json, name=risposta]
{
  "data": {
    "professors": {
      "data": [
        { "attributes": { "lastName": "Starace" } },
        { "attributes": { "lastName": "Abelson" } }
      ]
    }
  }
}
\end{lstlisting}

Nessun campo superfluo: è la risposta all'over-fetching.
\end{esempio}

\begin{esempio}{Risorse collegate in una sola richiesta}
Cognome di ciascun professore \textbf{e} titolo dei suoi corsi:

\begin{lstlisting}[style=plain, name=query]
query {
  professors {
    data {
      attributes {
        lastName,
        courses {
          data {
            attributes {
              title
            }
          }
        }
      }
    }
  }
}
\end{lstlisting}

Una sola richiesta al posto di $1 + n$: è la risposta all'under-fetching.
\end{esempio}

\begin{esempio}{Filtri}
L'equivalente GraphQL di

\begin{lstlisting}[style=sql]
SELECT title, credits FROM courses
WHERE credits >= 6
  AND (title CONTAINS 'Web' OR title CONTAINS 'TypeScript')
\end{lstlisting}

è:

\begin{lstlisting}[style=plain, name=query]
query {
  courses(filters: {
    credits: { gte: 6 },
    or: [
      { title: { containsi: "Web" } },
      { title: { containsi: "TypeScript" } }
    ]
  }) {
    data {
      attributes {
        title, credits
      }
    }
  }
}
\end{lstlisting}
\end{esempio}

\section{Generazione automatica di siti statici}

A volte usare un CMS completo è eccessivo. In alcuni scenari i dati non cambiano
spesso: si pensi a un ristorante che aggiorna il menu sul proprio sito circa una
volta al mese.

Per un mese intero, ogni caricamento di pagina provoca un'interrogazione al
database, che restituisce sempre la stessa lista, dopodiché il CMS mostra
sempre lo stesso contenuto. Il CMS sta sprecando cicli di calcolo --- cioè
denaro ed emissioni --- per generare ripetutamente le stesse pagine. I
meccanismi di caching aiutano, ma restano un sovraccarico e la necessità di
tenere il CMS aggiornato.

Un approccio più efficiente è:

\begin{enumerate}
  \item generare le pagine \textbf{una volta sola}, quando i dati cambiano;
  \item servire agli utenti le pagine \textbf{statiche} già compilate.
\end{enumerate}

È la \term{generazione automatica di siti statici}.

\begin{figure}[H]
  \centering
  \begin{tikzpicture}[font=\Lato\scriptsize,
      n/.style={wtnode, minimum width=2.4cm, minimum height=1.0cm},
      f/.style={wtnodeplain, minimum width=2.2cm, minimum height=1.0cm},
      a/.style={-{Stealth[length=5pt,width=4pt]}, line width=0.85pt,
                draw=neutral!45!black}]

    \node[f] (c) at (0,0.6)  {contenuti};
    \node[f] (t) at (0,-0.7) {template};
    \node[n, draw=orange-gradient-6, fill=orange-gradient-1!40!white]
      (g) at (3.6,0) {\textbf{generatore}};
    \node[f] (s) at (7.2,0) {sito statico\\ \scriptsize HTML, CSS, JS};
    \node[n] (w) at (10.8,0) {\textbf{web server}};

    \draw[a] (c.east) -- (g.west |- c);
    \draw[a] (t.east) -- (g.west |- t);
    \draw[a] (g) -- node[wtlab, above]{build} (s);
    \draw[a] (s) -- node[wtlab, above]{deploy} (w);
  \end{tikzpicture}
  \caption{Il flusso di un generatore di siti statici.}
  \label{fig:ssg}
\end{figure}

Fra i generatori più noti: \term{Hugo} (scritto in Go), \term{Gatsby} (genera
siti statici con React), \term{Jekyll} (basato su Ruby), \term{Pelican}
(Python) e \term{Zola} (Rust).

\section{Riferimenti}

\begin{itemize}
  \item \textbf{Definition of CMS} --- MDN web docs. \\
        \urlx{https://developer.mozilla.org/en-US/docs/Glossary/CMS}
  \item \textbf{What is Headless CMS?} --- AWS docs. \\
        \urlx{https://aws.amazon.com/what-is/headless-cms/}
  \item \textbf{The Fullstack Tutorial for GraphQL}. \\
        \urlx{https://www.howtographql.com/} \\
        Parti rilevanti: \textit{GraphQL Fundamentals}.
  \item \textbf{What is a static site generator?} --- Cloudflare Learning. \\
        \urlx{https://www.cloudflare.com/learning/performance/static-site-generator/}
\end{itemize}
